\Needspace{32\baselineskip}
\section{Proof checking and remaining obligations}
\Needspace{24\baselineskip}
\subsection{The two axiom audits}
The file \leanfile{Hypernets/Audit.lean}{Audit.lean} contains six
\texttt{\#print axioms} commands. They inspect counting-simulation soundness,
the forward and converse simulation constructions, the two contextual sufficiency
theorems, and the zero-step converse regression theorem. These diagnostic commands
introduce no new definitions or proofs.
\lean{Audit}{5}{10}{Audit: abstract theorem dependencies}
The graph audit inspects shallow plugging, host and inserted path equivalence,
plugging congruence, exact trace uniqueness, typed operation-path uniqueness,
finite deep enumeration, unique holes under boxing, deep plugging and its other-label
locality, and the self-loop and deep-replacement regression proofs. Twelve further commands
audit occurrence-coordinate observation preservation, exact hole inventories,
sharp freshness and consumption, isomorphism invariance of contexts, deep site
provenance, deep context closure, and the new reuse and collision regressions.
\lean{Graph/Audit}{5}{17}{Graph.Audit: concrete theorem dependencies}
\lean{Graph/Audit}{19}{30}{Graph.Audit: context calculus dependencies}
Nine further commands audit the exact deep inventory, both versions of sharp
context closure, hole-freeness, selected-label consumption, certified deep
plugging, and the deep reuse and collision regressions.
\lean{Graph/Audit}{32}{40}{Graph.Audit: exact deep inventory dependencies}
Eleven further commands audit recursive label transitivity, operation-path
invariance, occurrence embeddings, hole and context invariance, deep congruence,
quotient plugging and context validity, and the strict-extension regressions.
\lean{Graph/Audit}{42}{52}{Graph.Audit: recursive isomorphism dependencies}
Ten further commands audit shallow host congruence, site transport and its
specification, simultaneous congruence, quotient invariance, retained holes,
freshness, result validity, and the moved-address and wrong-address regressions.
\lean{Graph/Audit}{54}{63}{Graph.Audit: site transport dependencies}
Sixteen further commands inspect shallow associativity, both corolla units,
deep associativity, quotient composition, composed retained inventories, sharp
composite context closure and its certified constructor, three-stage coherence,
deep restoration, quotient identity, and the self-loop, repair, and mixed-depth
regressions.
\lean{Graph/Audit}{65}{80}{Graph.Audit: composition and unit dependencies}

Six further commands inspect global focus-free substitution, focussed shallow
and deep substitution, finite-filling boundary agreement, the initiality
regression, and rejection of a boxed focus.
\lean{Graph/Audit}{82}{87}{Graph.Audit: concrete focus dependencies}

Eight further commands inspect operation-view existence, search determinism,
search boundary laws, shallow filling preservation, and the thunk, value, and
variable regressions.
\lean{Graph/Audit}{89}{96}{Graph.Audit: search rule dependencies}

Ten further commands inspect explicit focus insertion and certification, both
incidence-recovery laws, hole freedom, term-typing preservation, initial
attachment, search endpoint realization, and the interior and value regressions.
\lean{Graph/Audit}{98}{107}{Graph.Audit: focus realization dependencies}

Twelve further commands inspect general erasure and its invariants, the
reinsertion isomorphism and focus endpoints, arbitrary-state recovery, state
search determinism and boundary laws, and the execution/interior regressions.
\lean{Graph/Audit}{109}{120}{Graph.Audit: erasure and state search dependencies}

Fifteen further commands inspect insertion--erasure cancellation, coordinate
transport and realization, erasure naturality, source invariance, quotient
search determinism and boundaries, and the concrete quotient regressions.
\lean{Graph/Audit}{122}{136}{Graph.Audit: round trips and quotient search dependencies}

Fourteen further commands inspect recursive focus freedom, coordinate search,
insertion and erasure, recursive state search and quotient determinism, the
quotient projection and lifting laws, exact-length executions, and regressions.
\lean{Graph/Audit}{138}{151}{Graph.Audit: recursive state search and execution dependencies}

Sixteen further commands inspect canonical roots, search preservation of
rootedness, operation reachability, insertion of rooted paths, the concrete
path obstruction, refocusing, and the positive and negative state examples.
\lean{Graph/Audit}{153}{168}{Graph.Audit: rooted-state and path-obstruction dependencies}

In the pinned snapshot, counting-simulation soundness, finite deep enumeration,
unique holes under boxing, occurrence-coordinate observation preservation,
isomorphism invariance of unique holes, and uniqueness of the two-hole example
report only \texttt{propext} and \texttt{Quot.sound}.
Recursive label transitivity, recursive operation-path invariance, the
nonexistence of a literal-label isomorphism between the example boxes, and their
equality as quotient classes likewise report only these two axioms. Existence
of transported sites also reports only these two axioms; the chosen transport
and substitution results additionally use classical choice.
All audited composition and unit results report these three standard
axioms. The focus results do too, except the boxed-focus rejection, which uses
only \texttt{propext} and \texttt{Quot.sound}. The search audits use only \texttt{propext} and \texttt{Quot.sound}, except
shallow filling preservation, which also uses classical choice. The new realization audits use all three standard axioms, except term-typing
preservation, which uses only propositional extensionality and quotient soundness.
All new general-erasure and state-search audit entries use those three
standard axioms. The round-trip and quotient-search audits do too, except
coordinate search renaming, which uses only \texttt{propext} and
\texttt{Quot.sound}. The remaining audited declarations additionally report \texttt{Classical.choice}.
The recursive-search audits use all three axioms, except label and graph
focus-freedom invariance, coordinate search transport, and the renamed-box
focus-freedom example, which use only \texttt{propext} and \texttt{Quot.sound}.
The rooted-state audit entries use all three axioms, except coordinate
operation-reachability preservation and the incoming-trace operation theorem,
which use only \texttt{propext} and \texttt{Quot.sound}.
These are standard Lean logical axioms: propositional extensionality, quotient
soundness, and classical choice, respectively. Neither audit reports
\texttt{sorryAx} or a project-specific axiom.

Absence of project-specific axioms does not discharge a theorem's hypotheses.
Specimen laws and initial lifting remain parameters of the abstract sufficiency
results. They do not appear as named global axioms in the report. In contrast,
the concrete graph constructions explicitly prove the invariant fields needed
to construct their results. Neither fact supplies the missing bridge to a
concrete UAM.

\Needspace{24\baselineskip}
\subsection{Reproducing the Lean and document checks}
The snapshot pins Lean \texttt{4.35.0} in
\leanfile{lean-toolchain}{lean-toolchain} and uses only the bundled \texttt{Std}.
With that toolchain available, run
\begin{quote}\ttfamily
lake build\\
lake env lean Hypernets/Audit.lean\\
lake env lean Hypernets/Graph/Audit.lean
\end{quote}
The script \leanfile{scripts/check.sh}{scripts/check.sh} runs the build and both
audits. Run it using \texttt{bash} with the script path as its argument. It selects
the workspace-local
macOS ARM compiler if present, and otherwise uses \texttt{lake} from the environment.
The audits are checked separately from the root import.

The commentary's source checker recursively scans all Lean modules under
\texttt{Hypernets/}, including nested graph modules, attributed declarations,
noncomputable definitions, abbreviations, and anonymous examples. It verifies
that each declaration has a source-fragment link, all linked ranges are valid,
and each working source file equals the version at the pinned commit. Diagnostic
audit ranges are linked separately. Run it with
\begin{quote}\ttfamily\small
python3 docs/mathematical-commentary/check\_sources.py
\end{quote}
This syntactic coverage check complements the mathematical exposition; it does
not itself check the correctness of English explanations. Source links are
immutable GitHub commit links, so future edits cannot move the referenced proof.

From the commentary project directory, \texttt{make} checks the links, compiles
with \texttt{latexmk}, and exports the PDF to the repository's
\texttt{output/pdf/} directory. The committed PDF can therefore be read without
a TeX installation, and its editable sources and reproduction instructions are
included alongside it.

\Needspace{24\baselineskip}
\subsection{What remains for the paper's concrete theorem}
This is a to-do list only: completed obligations are removed. The remaining
work needed to claim Theorem~6 with its original hypotheses is:
\begin{enumerate}
\item Prove representation adequacy between the typed depth-indexed hierarchy
with explicit binder metadata and the paper's set-based graph-labelled
hypernets, including compatibility across different depth bounds.
\item Construct a quotient of host/site pairs and prove interchange of
substitutions at independent sites.
\item Identify erasure-based state search with the paper's contextual local
graph rewrites in both directions.
\item Define the concrete substitution and compute transitions, with their
typing and context restrictions, and assemble the full UAM transition relation.
\item Prove rootedness preservation for substitution and the required
full-UAM structural consequences of determinism and refocusing. Establish the
one-way and output-closed specimen cases of Lemma~26(2--3), and the stationary
rewrite criterion where needed.
\item Define concrete pre-templates, specimens, quasi-specimens, output-closure,
input-safety, and robustness with the paper's restrictions.
\item Construct the abstract specimen interface from these objects. Connect
concrete specimen substitution to the finite-filling boundary theorem for
Lemma~39(1), and prove the rootedness and context-only traversal laws needed
for Lemma~39(2--3), including entering/exiting classifications. Prove the safety
and robustness bridges and the concrete observation/context membership
witnesses needed for initial lifting.
\item Instantiate the checked contextual sufficiency results using the
concrete interface, observations, and lifting. Discharge every remaining
concrete obligation to obtain the target robustness theorem.
\end{enumerate}
